\chapter{Croissance de la norme des coefficients}
\label{ann:norme}

\begin{lecture}
Le plafond d'itérations du chapitre~\ref{ch:entrainer} paraît arbitraire tant
qu'on ignore le cas qu'il traite. La norme des coefficients cesse de croître par un équilibre entre deux effets, et
cet équilibre disparaît lorsque les deux groupes sont parfaitement séparables.
\end{lecture}

\section{Équilibre qui arrête la croissance}

Les coefficients grandissent à chaque tour : leur norme passe de $0{,}30$ au
premier tour à $4{,}89$ au quatre centième, alors que la frontière est en place
depuis longtemps et que plus aucune décision ne change. Pourquoi ne
grandissent-ils pas indéfiniment ?

La réponse commande un morceau du programme. Tant qu'il reste des élèves mal
classés, la croissance s'arrête d'elle-même et le critère de stagnation termine
la boucle. Lorsque tous sont bien classés, elle se poursuit indéfiniment et seul
le plafond d'itérations arrête la descente : c'est la raison pour laquelle le
code de la deuxième partie porte ces deux conditions d'arrêt plutôt qu'une
seule.

Agrandir les coefficients raidit la sigmoïde autour de la frontière, sans
déplacer celle-ci. Les probabilités s'écartent donc de $0{,}5$, ce qui produit
deux effets opposés : les $1\,564$ élèves bien classés voient la leur tendre vers
$1$ et leur coût vers $0$ ; les $36$ autres voient la leur tendre vers $0$ et leur
coût vers l'infini.

\begin{figure}[htbp]
\centering
\begin{tikzpicture}[font=\footnotesize,y=2.3cm]
\draw[ink!35] (-0.4,0) -- (9.6,0);
\node[anchor=east,font=\scriptsize,text=ink!70] at (-0.45,0) {coût $0$};
\foreach \x/\h/\lab in {0.3/0.233/{$0{,}233$}, 1.2/0.067/{$0{,}067$},
                        2.1/0.018/{$0{,}018$}}{
  \fill[accent!70] (\x,0) rectangle (\x+0.6,\h);
  \node[anchor=south,text=accent,font=\scriptsize] at (\x+0.3,\h+0.02) {\lab};
}
\foreach \x/\h/\lab in {5.6/0.851/{$0{,}851$}, 6.5/1.029/{$1{,}029$},
                        7.4/1.228/{$1{,}228$}}{
  \fill[warning!70] (\x,0) rectangle (\x+0.6,\h);
  \node[anchor=south,text=warning,font=\scriptsize] at (\x+0.3,\h+0.02) {\lab};
}
\foreach \x/\m in {0.3/1, 1.2/2, 2.1/3, 5.6/1, 6.5/2, 7.4/3}{
  \node[anchor=north,font=\scriptsize,text=ink!70] at (\x+0.3,-0.02)
    {norme $\times\m$};
}
\node[anchor=north,align=center,text=ink] at (1.2,-0.16)
  {élève $3$, \Gryf, bien classé\\ \scriptsize son coût tend vers $0$};
\node[anchor=north,align=center,text=ink] at (6.5,-0.16)
  {élève $1239$, mal classé\\ \scriptsize son coût croît sans limite};
\draw[ink!30,dashed] (2.9,-0.05) -- (2.9,1.35);
\end{tikzpicture}
\caption{Le coût de deux élèves lorsque la norme des coefficients est multipliée
par $1$, $2$ puis $3$, la frontière restant fixe. À gauche, un élève du bon côté :
son coût se rapproche de zéro et ne peut pas descendre plus bas, si bien que tout
le gain possible tient dans les $0{,}233$ de départ. À droite, un élève du
mauvais côté : le sien augmente à chaque fois, et rien ne le borne. Multiplier
encore par dix porterait le premier à $0{,}000$ et le second à $2{,}99$.}
\label{fig:deux-eleves}
\end{figure}
\FloatBarrier

\begin{figure}[htbp]
\centering
\begin{tikzpicture}
\begin{axis}[width=.82\textwidth,height=7cm,
  axis lines=left,axis line style={ink!40},
  grid=major,grid style={gridgray!30},
  tick label style={font=\small},label style={font=\small},
  xlabel={norme des coefficients $\|w\|$},ylabel={coût},
  xmin=0,xmax=14.7,ymin=0,ymax=0.30,
  legend style={at={(0.98,0.97)},anchor=north east,draw=none,font=\small,
    cells={anchor=west},fill=white,fill opacity=0.85,text opacity=1}]
  \addplot[ink,very thick,no markers] table[x=norme,y=total] {data/echelle.dat};
  \addlegendentry{coût total}
  \addplot[accent,thick,densely dashed,no markers]
    table[x=norme,y=right] {data/echelle.dat};
  \addlegendentry{part des $1\,564$ élèves bien classés}
  \addplot[warning,thick,densely dotted,no markers]
    table[x=norme,y=wrong] {data/echelle.dat};
  \addlegendentry{part des $36$ élèves mal classés}
  \addplot[only marks,mark=*,ink,mark size=2.2pt] coordinates {(5.384,0.1073)};
  \draw[ink!55,thin,-{Latex[length=1.6mm]}]
    (axis cs:8.2,0.055) -- (axis cs:5.7,0.101);
  \node[font=\scriptsize,text=ink,anchor=west] at (axis cs:8.3,0.052)
    {minimum en $\|w\|=5{,}38$};
\end{axis}
\end{tikzpicture}
\caption{Les deux parts du coût lorsque les coefficients sont agrandis, la
frontière restant fixe. Celle des élèves bien classés décroît vers zéro sans
jamais passer dessous ; celle des $36$ autres croît sans borne, puisque
$-\log p$ diverge quand leur probabilité tend vers $0$. La somme atteint son
minimum en $\|w\|=5{,}38$.}
\label{fig:deux-parts}
\end{figure}
\FloatBarrier

\begin{table}[htbp]
\centering
\begin{tabular}{cccc l}
\toprule
$\|w\|$ & $1\,564$ bien classés & $36$ mal classés & total & \\
\midrule
$0{,}49$  & $0{,}445$ & $0{,}019$ & $0{,}464$ & transition très étalée \\
$4{,}89$  & $0{,}041$ & $0{,}067$ & $0{,}108$ & modèle de ce chapitre \\
$5{,}38$  & $0{,}034$ & $0{,}073$ & $0{,}107$ & minimum \\
$14{,}68$ & $0{,}005$ & $0{,}189$ & $0{,}194$ & transition abrupte \\
\bottomrule
\end{tabular}
\caption{Entre la première ligne et la dernière, la part des élèves bien classés
est divisée par quatre-vingt-neuf, celle des mal classés multipliée par dix.}
\label{tab:deux-parts}
\end{table}

Le gain est borné par zéro, la perte ne l'est pas. Au-delà de $\|w\|=5{,}38$,
agrandir les coefficients coûte donc plus qu'il ne rapporte, et la descente cesse
de le faire. La descente l'atteint en suivant le gradient. Le
modèle de ce chapitre s'arrête un peu avant, à $4{,}89$, parce que le compteur
est coupé au quatre centième tour.

Ce mécanisme disparaît lorsqu'une droite sépare parfaitement les deux groupes ---
les données sont alors dites \emph{linéairement séparables}. La part des mal
classés est vide, la courbe orange de la figure~\ref{fig:deux-parts} est
identiquement nulle, et il ne reste que la courbe bleue, décroissante : le coût
diminue indéfiniment, $\|w\|$ diverge, et seul le plafond d'itérations arrête la
descente (ch.~\ref{ch:convexite}).

\begin{vigilance}[Reconnaître et traiter la divergence]
Ce cas se diagnostique à trois signes conjoints : la descente consomme tout son
plafond d'itérations sans jamais déclencher son critère d'arrêt, les coefficients
continuent de croître régulièrement, et les probabilités affichées valent
exactement $0$ ou $1$ après arrondi. La cause est alors dans les données, pas
dans le programme.

Deux remèdes existent. Le plus simple est d'arrêter sur un nombre de tours fixé,
ce que fait le plafond : la frontière obtenue est correcte, seule sa raideur est
arbitraire. Le second consiste à ajouter au coût un terme $\frac{\mu}{2}\|w\|^2$
qui croît avec la norme --- c'est la \emph{régularisation}. La courbe bleue de la
figure~\ref{fig:deux-parts} est alors relevée par une parabole, un minimum
réapparaît même sur des données séparables, et le paramètre $\mu$ fixe où il
tombe. La bibliothèque \texttt{scikit-learn} l'applique par défaut, ce qui
explique qu'elle converge sur des données où une implémentation directe diverge.
\end{vigilance}

\begin{resultat}[Interprétation de la norme finale]
La norme finale traduit le degré de séparation que les données justifient. Un
recouvrement des deux populations la maintient basse et laisse une transition
étalée ; une séparation nette la porte haut. La forcer davantage ferait annoncer
une certitude que les données ne soutiennent pas, et une probabilité de $0{,}99$
sur un cas limite est une information fausse.
\end{resultat}

